\documentclass[11pt]{article}

\usepackage[margin=1in]{geometry}
\usepackage[T1]{fontenc}
\usepackage{lmodern}
\usepackage{xcolor}
\usepackage{sectsty}
\usepackage{hyperref}
\usepackage{enumitem}
\usepackage{parskip}
\usepackage{titling}

\definecolor{headingblue}{RGB}{40,85,140}
\definecolor{alertred}{RGB}{180,40,40}
\definecolor{codegray}{RGB}{240,240,240}

\sectionfont{\color{headingblue}\normalsize\bfseries}
\subsectionfont{\color{headingblue}\small\bfseries}

\hypersetup{
  colorlinks=true,
  urlcolor=headingblue,
  linkcolor=headingblue
}

\setlength{\parindent}{0pt}
\setlength{\parskip}{0.6em}

\newcommand{\cmd}[1]{\par\vspace{0.2em}\colorbox{codegray}{\parbox{\dimexpr\linewidth-2\fboxsep}{\ttfamily\small #1}}\par\vspace{0.2em}}
\newcommand{\alert}[1]{{\color{alertred}\textbf{\rule[-0.2ex]{0.8ex}{2ex}\ #1}}}
\newcommand{\plat}[1]{{\color{headingblue}\textbf{#1}}}

\title{\vspace{-2em}\Huge\bfseries Session 4 Pre-Work\\[0.3em]\Large\mdseries\itshape Skills, MCPs, and the Collaborative Story}
\author{}
\date{}

\begin{document}

\maketitle
\vspace{-2em}
\begin{center}
Instructor: Brady Allardice --- \href{mailto:brady.allardice@upf.edu}{brady.allardice@upf.edu}
\end{center}

\hrule
\vspace{0.8em}

Session 4 introduces \textbf{MCPs} --- the way Claude Code talks to systems outside your project. We will use two of them live: \textbf{GitHub} (for a collaborative writing exercise) and \textbf{Zotero} (for a citations demo). Both need to be installed and authenticated before class --- they take 15--20 minutes each and the setup is too fiddly to do in the room.

Please complete \textbf{all six steps} below. Budget \textbf{45--60 minutes total}. If you get stuck for more than 10 minutes, stop and email me --- don't try to debug it yourself. Once you are done, run the verification commands at the end and \textbf{send me a screenshot} of the output \emph{and} your GitHub username (Step 2).

\section*{What You Need Before Starting}
\begin{itemize}[leftmargin=2em,itemsep=0pt]
  \item Everything from the original Pre-Workshop Setup Guide working (Claude Code authenticated, VS Code running, Python installed)
  \item A GitHub account (created during initial setup)
  \item A free Zotero account --- we'll create one in Step 5 if you don't have one
\end{itemize}

\alert{If \texttt{claude --version} or \texttt{git --version} doesn't work, fix that first --- these steps assume the original setup is intact.}

\section*{Step 1: Refresh the Course Repo and Set Up Your Workspace}

Starting with Session 4 we are switching to a new workflow: the course repo stays \textbf{read-only} on your side, and you do all your actual work in a separate \texttt{\textasciitilde/my\_workspace/} folder. This means I can push new material any time without breaking whatever you have in progress, and you can edit/break/delete freely without worrying about merge conflicts against the shared repo.

\textbf{Run these three commands now}, and plan to run them at the start of every future session too:

\cmd{\# 1. Pull the latest course repo\\cd \textasciitilde/claude-code-workshop \&\& git pull\\\\\# 2. Copy today's session folder into your workspace\\mkdir -p \textasciitilde/my\_workspace \&\& cp -r session\_4/ \textasciitilde/my\_workspace/session\_4/\\\\\# 3. Open your copy in VS Code and work there\\cd \textasciitilde/my\_workspace/session\_4 \&\& code .}

\textit{If you cloned the repo somewhere other than \texttt{\textasciitilde/claude-code-workshop} (e.g.\ inside Dropbox or Documents), substitute your actual path in step 1.} If you don't already have the course repo cloned at all, run this first:
\cmd{cd \textasciitilde{} \&\& git clone https://github.com/bradyallardice/claude-code-workshop.git}

After Step 2, you should see a \texttt{session\_4/} folder inside \texttt{\textasciitilde/my\_workspace/} containing the new material. After Step 3, VS Code should open in that copy. From now on, \textbf{edit files in \texttt{\textasciitilde/my\_workspace/session\_4/}, not in the course repo itself.}

If \texttt{git pull} fails with a merge conflict, that is exactly the problem we fix in class --- don't try to resolve it. Just email me and I will help you reset cleanly.

\section*{Step 2: Send Me Your GitHub Username}

In Session 4 we run a \textbf{collaborative writing exercise} using a separate class GitHub repository. I need to add each of you as a collaborator before class so you can push branches and open pull requests.

\textbf{Email me your GitHub username} (not your email --- the username, e.g.\ \texttt{octocat}) at \href{mailto:brady.allardice@upf.edu}{brady.allardice@upf.edu}. You can find it in the top-right of \textbf{https://github.com} after logging in.

\section*{Step 3: Update Claude Code}

The MCP install commands below assume Claude Code 2.1 or newer. Update to the latest version:
\cmd{claude update}

Then confirm the version:
\cmd{claude -\/-version}

\textit{Expected: 2.1.x or higher.}

\section*{Step 4: Install the GitHub MCP}

The GitHub MCP lets Claude Code open pull requests, leave review comments, and create issues directly from the chat. We use it during the collaborative story exercise.

\subsection*{4a. Create a GitHub Personal Access Token}

A Personal Access Token (PAT) is a password-like string that lets the MCP act on your behalf. \textbf{Treat it like a password} --- don't paste it into chat, don't commit it to git.

\begin{itemize}[leftmargin=2em,itemsep=0pt]
  \item Go to \textbf{https://github.com/settings/tokens?type=beta}
  \item Click \textbf{Generate new token}
  \item \textbf{Token name:} \texttt{claude-code-workshop}
  \item \textbf{Expiration:} 90 days
  \item \textbf{Repository access:} ``All repositories'' is fine for the workshop
  \item \textbf{Repository permissions:} grant \textbf{Read and write} on \emph{Contents}, \emph{Issues}, and \emph{Pull requests}
  \item Click \textbf{Generate token}, then \textbf{copy the token immediately} --- GitHub will not show it again
\end{itemize}

Paste it into a temporary scratch file you will delete after Step 4b.

\subsection*{4b. Add the GitHub MCP to Claude Code}

In your terminal, paste the command below, \textbf{but before pressing Enter, replace the literal text \texttt{YOUR\_PAT} with the token you just copied} (keep the word \texttt{Bearer} and the space before the token --- you are only replacing \texttt{YOUR\_PAT}, nothing else):
\cmd{claude mcp add -\/-transport http github https://api.githubcopilot.com/mcp/ -H "Authorization: Bearer YOUR\_PAT"}

\textit{Example of what it should look like after substitution (with a fake token):}
\cmd{claude mcp add -\/-transport http github https://api.githubcopilot.com/mcp/ -H "Authorization: Bearer github\_pat\_11ABCDEFG0abc123XYZ..."}

Then verify:
\cmd{claude mcp list}

You should see \texttt{github} in the list with a \texttt{$\checkmark$ Connected} status. \textbf{Now delete the scratch file containing your token.}

\subsection*{4c. Confirm it works}
Open Claude Code in any directory and ask:
\cmd{Use the GitHub MCP to list the five most recent public repos owned by my account.}

If it returns a list, you are done. If it asks you to authenticate or returns an error, see Troubleshooting at the end.

\section*{Step 5: Install the Zotero MCP}

The Zotero MCP lets Claude Code search your Zotero library, pull BibTeX, and insert citations into a draft. We use it for a citations demo and you will use it for your own papers afterward.

\subsection*{5a. Set up Zotero and import the practice library}

\begin{itemize}[leftmargin=2em,itemsep=0pt]
  \item If you don't have a Zotero account, create a free one at \textbf{https://www.zotero.org/user/register}
  \item Install the \textbf{Zotero desktop app} from \textbf{https://www.zotero.org/download/} if you don't already have it --- the MCP needs the desktop app running to read your library (including full-text PDFs)
  \item Open the desktop app and sign in
\end{itemize}

\textbf{Now import the practice library.} I have prepared a bundle of 10 papers (with PDFs and notes) for the in-class demo. It lives in the course repo at:

\cmd{session\_4/docs/claude\_code\_zotero/}

After running \texttt{git pull} from Step 1, that folder contains a \texttt{claude\_code\_zotero.rdf} file and a \texttt{files/} subfolder of PDFs.

\textbf{To import:}
\begin{itemize}[leftmargin=2em,itemsep=0pt]
  \item In the Zotero desktop app, click \textbf{File $\rightarrow$ Import\ldots}
  \item Choose \textbf{``A file (BibTeX, RIS, Zotero RDF, etc.)''}, click \textbf{Continue}
  \item Navigate to \texttt{session\_4/docs/claude\_code\_zotero/} and select \texttt{claude\_code\_zotero.rdf}
  \item On the next screen, check \textbf{``Place imported collections and items into new collection''} and \textbf{``Copy files to the Zotero storage folder''}
  \item Click \textbf{Continue}
\end{itemize}

You should end up with a new collection called \texttt{claude\_code\_zotero} in the left sidebar containing 10 items, each with a PDF attached. \textbf{Confirm by clicking on the collection and verifying you see all 10 papers with attachment paperclip icons.}

\alert{Critical:} keep the \texttt{files/} folder \emph{next to} the \texttt{.rdf} on disk during import --- if you move just the .rdf or import it from elsewhere, Zotero will import the metadata but not the PDFs.

\subsection*{5b. Let Claude Code talk to your Zotero app}

Zotero blocks other applications from connecting by default. Turn that off:

\begin{itemize}[leftmargin=2em,itemsep=0pt]
  \item Open the Zotero desktop app
  \item \plat{Mac:} \textbf{Zotero $\rightarrow$ Settings} \quad \plat{Windows:} \textbf{Edit $\rightarrow$ Settings}
  \item Go to the \textbf{Advanced} tab
  \item Under \textbf{General}, check the box \textbf{``Allow other applications on this computer to communicate with Zotero''}
  \item Close the Settings window. \textbf{Leave Zotero running} --- the MCP needs it open to query your library
\end{itemize}

\alert{The MCP also falls back to the Zotero Web API if the desktop app is closed, but the local connection is faster and gives Claude Code access to your PDFs --- worth turning on.}

\subsection*{5c. Get your Zotero API key and Library ID}

\begin{itemize}[leftmargin=2em,itemsep=0pt]
  \item Go to \textbf{https://www.zotero.org/settings/keys}
  \item Your \textbf{Library ID} (a number, e.g.\ \texttt{1234567}) is shown at the top under ``Your userID for use in API calls''
  \item Click \textbf{Create new private key}
  \item \textbf{Key name:} \texttt{claude-code-workshop}
  \item Check \textbf{Allow library access} (read-only is fine; check write access too if you want Claude to add items)
  \item Click \textbf{Save Key}, then \textbf{copy the key immediately}
\end{itemize}

Save both the key and your Library ID into a temporary scratch file.

\subsection*{5d. Install the \texttt{zotero-mcp} package}

\plat{Mac:}
\cmd{pip install zotero-mcp}
If \texttt{pip} is not found, use:
\cmd{pip3 install zotero-mcp}

\plat{Windows:}
\cmd{py -m pip install zotero-mcp}

Then confirm it installed:
\cmd{zotero-mcp -\/-help}

\textit{Expected: a help message listing options. If you see ``command not found,'' close and reopen your terminal first.}

\subsection*{5e. Add the Zotero MCP to Claude Code}

Run (substituting your real values from Step 5c):
\cmd{claude mcp add zotero -e ZOTERO\_LIBRARY\_ID=YOUR\_LIBRARY\_ID -e ZOTERO\_API\_KEY=YOUR\_API\_KEY -e ZOTERO\_LIBRARY\_TYPE=user -- zotero-mcp}

Then verify:
\cmd{claude mcp list}

You should now see both \texttt{github} and \texttt{zotero} listed with \texttt{$\checkmark$ Connected}. \textbf{Delete the scratch file with your key and Library ID.}

\subsection*{5f. Confirm it works}

Open Claude Code in any directory and ask:
\cmd{Use the Zotero MCP to list 5 items from the claude\_code\_zotero collection.}

If it returns titles from the practice library (Alt et al., Bisgaard, Hicks et al., etc.), you are done. \textbf{Make sure the Zotero desktop app is open} when you run this --- if it isn't, the MCP will either fail or fall back to the (slower) Web API.

\section*{Step 6: Re-Read the Session 3 Git Cycle (5 min)}

Session 4 leans heavily on the Git workflow from Session 3. Spend five minutes re-reading the slides from Session 3 (\texttt{slides/session3\_slides.pdf} in the course repo), specifically the \texttt{commit $\rightarrow$ instruct $\rightarrow$ diff $\rightarrow$ decide} cycle and how branches work. The collaborative story exercise will move faster than this if those concepts are fresh.

\section*{Verification: Confirm Everything Works}

Open a terminal and run each of the following. \textbf{Take a screenshot of the full output and send it to me} along with your GitHub username from Step 2.

\cmd{claude -\/-version}
\textit{Expected: \texttt{2.1.x} or higher.}

\cmd{claude mcp list}
\textit{Expected: both \texttt{github} and \texttt{zotero} listed with \texttt{$\checkmark$ Connected}.}

\cmd{zotero-mcp -\/-help}
\textit{Expected: a help message (any output is fine; ``command not found'' is the failure mode).}

\textbf{Also confirm:}
\begin{itemize}[leftmargin=2em,itemsep=0pt]
  \item \texttt{git pull} on the course repo succeeds and you see a \texttt{session\_4/} folder with new content
  \item The \texttt{claude\_code\_zotero} collection in your Zotero app contains 10 items, each with a PDF attached
  \item Asking Claude Code ``list 5 items from the claude\_code\_zotero collection'' returns titles
  \item Asking Claude Code ``list my five most recent GitHub repos'' returns repo names
  \item You have emailed me your GitHub username
\end{itemize}

\section*{Troubleshooting}

\textbf{\texttt{claude mcp add} fails with ``unknown flag'':} You are on an older Claude Code. Run \texttt{claude update} (Step 3) and try again.

\textbf{GitHub MCP shows ``Connected'' but tool calls fail with 401/403:} Your PAT is missing a permission. Regenerate it (Step 4a) and make sure \emph{Contents}, \emph{Issues}, and \emph{Pull requests} are all set to \emph{Read and write}. Then re-run the \texttt{claude mcp add} command --- it will overwrite the existing entry.

\textbf{Zotero MCP shows ``Connected'' but returns empty results:} Your library is genuinely empty, or \texttt{ZOTERO\_LIBRARY\_TYPE} is wrong. If you set up a \emph{group} library instead of a personal one, change \texttt{user} to \texttt{group} in the env var.

\textbf{\texttt{zotero-mcp: command not found} after \texttt{pip install}:} Close and reopen your terminal. If still failing, the install went to a directory not on your PATH --- email me with the output of \texttt{pip show zotero-mcp} and I will help.

\textbf{You accidentally pasted a token into a chat or committed it to git:} Revoke it immediately (GitHub: \textbf{https://github.com/settings/tokens}; Zotero: \textbf{https://www.zotero.org/settings/keys}) and generate a new one. Tokens are easy to rotate; a leaked token in a public repo is not.

\textbf{Anything else:} Email me at \href{mailto:brady.allardice@upf.edu}{brady.allardice@upf.edu}. Do not spend more than 10 minutes troubleshooting on your own. I would much rather help you before the workshop than lose time during it.

\end{document}
